% Focus: State the revised goals and contributions with precise scope; emphasize what is proved vs. deferred.

\subsection{Motivation and guiding examples}\label{subsec:intro-motivation}
% Focus: Reframe AEG as geometry of evaluation histories; give one concrete threadlike example and its two propagations.

Arithmetic expressions come with a hidden combinatorial choice: an \emph{evaluation history}. Even when two evaluation histories give the same final value, intermediate values can differ, and commutators of elementary steps can carry systematic defects. Arithmetic Expression Geometry (AEG) is a proposal to treat such defects as geometric data.

The present paper focuses on a distinguished class of \emph{threadlike} expressions, where evaluation histories naturally decompose into two elementary propagations: additive updates and multiplicative updates. Passing from discrete propagation to a continuous limit yields an arithmetic flow equation. The central theme is that the non-commutativity between these two propagations admits both global invariants and local differential-geometric incarnations.

\subsection{Main results and structure of the paper}\label{subsec:intro-results}
% Focus: List theorem-level contributions (flow, AES models, ACS identity, contact/horizontal calculus, horizontal holomorphic fields/twisted harmonicity, invariance).

The paper establishes a first foundational layer of AEG. Its contributions are organized into the following theorem-level packages.

\begin{itemize}
    \item \emph{Arithmetic flow equation.} From infinitesimal additive and multiplicative propagations we derive a first-order propagation law and reformulate it in a Lie-theoretic \texorpdfstring{$\mathrm{Aff}(1)$}{Aff(1)} framework.
    \item \emph{Arithmetic Expression Spaces (AES).} We define an AES as geometric data supporting an assignment field governed by the flow/eikonal law, and we construct the basic model spaces \texorpdfstring{$\mathfrak{E}_0$}{E0} and \texorpdfstring{$\mathfrak{E}_1$}{E1}, including verifiable computations in the upper half-plane.
    \item \emph{Arithmetic torsion and the ACS triple identity.} We define arithmetic torsion as a commutator defect of additive and multiplicative propagation. We introduce the Accumulative Commutative Space (ACS) and prove a Stokes-type triple identity expressing torsion equivalently as an evaluation difference, a weighted area integral, and a boundary integral under explicit hypotheses.
    \item \emph{Contact structure and horizontal calculus.} We lift the propagation data to a contact manifold whose horizontal distribution recovers the flow directions. This yields a horizontal differential calculus \texorpdfstring{$\delta$}{delta} and a local commutator curvature package reflecting the same arithmetic non-commutativity.
    \item \emph{Horizontal holomorphic fields and twisted harmonicity.} On the horizontal distribution we define an arithmetic Cauchy--Riemann system and package it into first-order operators. We show that this first-order system forces a nontrivial twisted second-order equation, providing an intrinsic analytic operator layer associated to arithmetic propagation.
    \item \emph{Normalization and invariance.} We systematize the effects of rescaling, change of logarithmic base, and coordinate changes, and we record the invariance properties of the core constructions and identities.
\end{itemize}

The remainder of the paper follows this structure. Sections~\ref{sec:paths-flow}--\ref{sec:aes} develop the propagation law and the definition and construction of AES. Section~\ref{sec:acs-torsion} develops arithmetic torsion and the ACS triple identity. Section~\ref{sec:contact-delta} develops the contact-geometric lift and the horizontal calculus. Section~\ref{sec:horizontal-holomorphic} introduces horizontal holomorphic fields and their induced twisted harmonicity. Section~\ref{sec:invariance-positioning} collects normalization, invariance, and mathematical positioning. Technical computations and examples are placed in the appendices.

\subsection{What is established here and what is deferred}\label{subsec:intro-scope}
% Focus: Explicitly separate proved statements from programmatic outlook (classification, full analysis, loop/relations theory).

This paper is foundational and intentionally delimits its scope.

\begin{itemize}
    \item We give explicit definitions, model constructions, and complete proofs for the central identities claimed in the main text.
    \item We do not attempt a global classification of AES beyond the model settings treated here. Classification and uniqueness problems are stated explicitly as open directions.
    \item The horizontal holomorphic and twisted harmonic operators introduced here are not presented as a full replacement of classical complex analysis. In particular, boundary integral kernels, analytic continuation, and global potential theory are deferred to later work.
    \item The multi-layered equality program and the associated loop/relations theory are only motivated by examples in this paper; a systematic global theory is deferred.
\end{itemize}

\subsection{Relation to existing frameworks}\label{subsec:intro-related}
% Focus: Briefly position relative to contact/jet spaces, CR/sub-Riemannian viewpoints, and rewriting/groupoid semantics.

Several parts of the construction naturally interface with established geometric structures. The contact form and horizontal calculus invite comparison with jet-space contact geometry and with CR/sub-Riemannian analysis on non-integrable distributions \cite{Darboux1882Pfaff,Geiges2001BriefHistory,Geiges2008IntroContact}. From a combinatorial viewpoint, evaluation histories and their commutators suggest a groupoid language for rewriting systems and a holonomy-type organization of equality levels \cite{Post1943FormalRO,Chomsky1956ThreeMF}. Section~\ref{sec:invariance-positioning} records these comparisons at the level needed for the present foundational development.
